\chapter{Les outils colorimétriques}

\section{RGB linéaire et changements de primaires}
\IPR{} travaille sur un master en RGB linéaire Rec.2020, D65 : les valeurs sont
proportionnelles à la lumière, sans courbe de transfert. Passer d'un jeu de
primaires à un autre est alors une simple application linéaire,
\begin{equation}
\mathbf{p}_{\text{sRGB}} = M\,\mathbf{p}_{\text{Rec.2020}},\qquad
M = M_{\text{sRGB}\to XYZ}^{-1}\,M_{\text{Rec.2020}\to XYZ},
\end{equation}
où chaque matrice $M_{\cdot\to XYZ}$ se déduit des coordonnées $xy$ des
primaires et du blanc~\cite{itu2020,iec61966}. Le passage n'écrête rien : une
couleur Rec.2020 hors du gamut sRGB obtient simplement une composante négative,
et l'aller-retour $M^{-1}M$ restitue la valeur exacte.

\resultat{Les tables de vibrance ont été mesurées sur des exports sRGB : \IPR{}
convertit le master en sRGB linéaire, applique le modèle, puis revient en
Rec.2020, sans écrêtage dans l'intervalle.}

\section{Pourquoi Oklab}
Pour agir sur « l'intensité d'une couleur » sans toucher à sa luminosité ni à
sa teinte, il faut un espace où ces trois notions sont séparées. Oklab,
proposé par Björn Ottosson en 2020~\cite{ottosson2020}, est construit pour que
\begin{itemize}
\item $L$ représente la clarté perçue ;
\item le plan $(a,b)$ porte la couleur : la distance au gris est la
\textbf{chroma} $C=\sqrt{a^2+b^2}$ et l'angle est la \textbf{teinte}
$h=\operatorname{atan2}(b,a)$ ;
\item une variation de $C$ à $L$ et $h$ fixes soit perçue comme une variation
d'intensité sans changement de teinte, mieux que dans HSV ou CIELAB.
\end{itemize}

Depuis le RGB linéaire sRGB, la transformation compte deux matrices et une
racine cubique :
\begin{align}
(l,m,s)^\top &= M_1\,(R,G,B)^\top,\\
(L,a,b)^\top &= M_2\,(\sqrt[3]{l},\sqrt[3]{m},\sqrt[3]{s})^\top,
\end{align}
avec
\begin{equation}
M_1=\begin{pmatrix}
0{,}4122 & 0{,}5363 & 0{,}0514\\
0{,}2119 & 0{,}6807 & 0{,}1074\\
0{,}0883 & 0{,}2817 & 0{,}6300\end{pmatrix},\quad
M_2=\begin{pmatrix}
0{,}2105 & 0{,}7936 & -0{,}0041\\
1{,}9780 & -2{,}4286 & 0{,}4506\\
0{,}0259 & 0{,}7828 & -0{,}8087\end{pmatrix}.
\end{equation}
\IPR{} utilise la racine cubique signée ($\sqrt[3]{-x}=-\sqrt[3]{x}$) : une
composante négative ou supérieure à 1, fréquente dans un master HDR, reste
calculable. Le retour inverse les deux étapes.

\section{Deux mesures d'intensité : chroma et saturation}
La \textbf{chroma} Oklab $C$ est une mesure absolue : elle grandit avec la
distance au gris. Une couleur sombre et une couleur claire de même aspect
« vif » n'ont pas la même chroma.

La \textbf{saturation} est une mesure relative. Le modèle utilise celle d'un
espace RGB à très large gamut, ProPhoto RGB (ROMM, ISO 22028-2~\cite{iso22028}) :
\begin{equation}
s=\frac{\max(R,G,B)-\min(R,G,B)}{\max(R,G,B)}\in[0,1],
\end{equation}
calculée en ProPhoto \emph{linéaire}. Le passage sRGB vers ProPhoto change de
blanc (D65 vers D50) par la transformation de Bradford~\cite{lam1985}.
$s=0$ désigne un gris, $s=1$ une couleur sur le bord du gamut ProPhoto.

Pourquoi ProPhoto ? Parce que c'est, parmi les variables testées, celle qui
explique le mieux le comportement positif de Lightroom (chapitre 3), ce qui
est cohérent avec un espace de travail interne fondé sur les primaires
ProPhoto. Ce n'est pas une affirmation sur le code de Lightroom, mais un
constat de mesure.

\resultat{Trois coordonnées décrivent un pixel pour la vibrance : sa
saturation ProPhoto $s$, sa clarté Oklab $L$ et sa teinte Oklab $h$. La chroma
$C$ sert à mesurer l'effet et à repérer les quasi-gris.}

\section{Rotation dans le plan de couleur}
Faire tourner la teinte d'un angle $\theta$ sans changer la chroma revient à
appliquer la matrice de rotation plane au couple $(a,b)$ :
\begin{equation}
\begin{pmatrix}a'\\ b'\end{pmatrix}=
\begin{pmatrix}\cos\theta & -\sin\theta\\ \sin\theta & \cos\theta\end{pmatrix}
\begin{pmatrix}a\\ b\end{pmatrix}.
\end{equation}
Cette matrice est orthogonale : elle conserve la norme, donc $C$, et ajoute
$\theta$ à $h$. Combinée à un gain $g$ sur la chroma et à un gain $k$ sur la
clarté, la transformation complète d'un pixel s'écrit
\begin{equation}
L'=k\,L,\qquad (a',b')^\top = g\,R(\theta)\,(a,b)^\top .
\label{eq:transformation}
\end{equation}
Gain et rotation commutent ($gR=Rg$) et restent indépendants : l'un règle
l'intensité, l'autre la teinte.

\resultat{Ce n'est pas une matrice unique appliquée à toute l'image : $g$, $k$
et $\theta$ dépendent du pixel et de la valeur du curseur. La transformation
globale n'est donc pas linéaire, même si elle l'est localement dans le plan
$(a,b)$.}
